\section{Model, clocks, and conventions}\label{sec:model}
This section fixes the normalization shared by the static and dynamical
arguments. In particular, adding a GOE diagonal will provide spectral
test vectors without changing any quantity intrinsic to the spin law.

Throughout, $n\to\infty$; arguments using three distinguished coordinates
are applied for $n\ge3$. Constants $c,C>0$ are independent of $n$, may
change from line to line, and have only the additional dependencies
explicitly indicated. Vector norms are Euclidean and matrix norms are
operator norms. We write $O_{\Prob}(1)$ for a sequence whose laws are
tight. Equivalently, its absolute value exceeds any positive deterministic
divergent threshold with probability tending to zero. The ability to
test every such threshold will be used in the covariance argument.

Following the Gaussian diagonal augmentation in
\cite[Equation~(1.1)]{comets1996}, let $W$ be a real symmetric GOE matrix
with independent entries on and above its diagonal, with
\[
 W_{ij}\sim N(0,1/n)\quad(i<j),\qquad
 W_{ii}\sim N(0,2/n).
\]
Set $\Omega_n=\{-1,1\}^n$ and
\[
 H_W(x)=\tfrac12x^TWx,
 \qquad
 \mu_W(x)=\frac{e^{H_W(x)}}{\sum_{y\in\Omega_n}e^{H_W(y)}}.
\]
We abbreviate $\mu_W$ to $\mu$ when the disorder is fixed. Because every $x_i^2=1$, the diagonal contributes the same quantity
$\frac12\sum_iW_{ii}$ to every configuration. It cancels from the Gibbs
normalization, leaving precisely \eqref{eq:intro-law}. Covariance,
heat-bath kernels and mixing times therefore depend only on the original
off-diagonal disorder. A spectral direction may still depend on the added
diagonal. In the realized-start proof we use that direction to test an
intrinsic total-variation distance, and only then integrate out the
diagonal. This distinction avoids imposing an unnecessary simultaneous
choice of a test direction for all initial states.

For functions $f,g:\Omega_n\to\R$, write
\[
 \langle f,g\rangle_\mu=\mu[fg],
 \qquad
 P_i f=\E_\mu[f\mid x_{-i}],
 \qquad
 \mathcal L=\sum_{i=1}^n(P_i-\mathrm{Id}),
 \quad P=\frac1n\sum_{i=1}^nP_i.
\]
The two transition operators are $e^{t\mathcal L}$ and $P^k$.
Their exact relation is $\mathcal L=n(P-\mathrm{Id})$; time in one clock
must therefore not be read as a number of attempts in the other.
Every conditional spin probability is positive. Single-site updates
connect the cube, and the discrete chain can hold, so both chains are
irreducible and the discrete chain is aperiodic.

\begin{lemma}[Heat-bath forms]\label{lem:heatbath-form}
Each $P_i$ is an orthogonal projection in $L^2(\mu)$. The unscaled Dirichlet form is
\begin{align}
 \mathcal D_\mu(f)
 &=-\langle f,\mathcal Lf\rangle_\mu
   =\sum_i\mu[(f-P_if)^2]  \label{eq:heatbath-form}\\
 &=\sum_i\inf_{h_i:\{-1,1\}^{n-1}\to\R}
       \mu[(f-h_i(x_{-i}))^2].\nonumber
\end{align}
The discrete form is $\langle f,(\mathrm{Id}-P)f\rangle_\mu=\mathcal D_\mu(f)/n$, and $P$ has spectrum in $[0,1]$. For a linear function,
\begin{equation}\label{eq:linear-form}
 \mathcal D_\mu(a^Tx)=\sum_i a_i^2\alpha_i,
 \qquad
 \alpha_i=\mu[\sech^2 g_i],\quad
 g_i(x)=\sum_{j\ne i}W_{ij}x_j.
\end{equation}
\end{lemma}
\begin{proof}
Conditional expectation is the orthogonal projection onto functions of
$x_{-i}$. Hence
$\langle f,(\mathrm{Id}-P_i)f\rangle_\mu=\|f-P_if\|_{L^2(\mu)}^2$,
and this squared distance is the infimum over all such functions.
Summing proves \eqref{eq:heatbath-form}; averaging instead proves the
discrete form identity. Moreover
$0\le\langle f,Pf\rangle_\mu\le\langle f,f\rangle_\mu$ for every $f$,
so the spectrum of $P$ lies in $[0,1]$.

For a linear observable, all terms except $a_ix_i$ are fixed after
conditioning on $x_{-i}$. The conditional mean and variance of $x_i$
are $\tanh g_i$ and $\sech^2g_i$, respectively. Its contribution to the
form is therefore $a_i^2\mu[\sech^2g_i]$, which gives
\eqref{eq:linear-form}.
\end{proof}

For probability laws $\nu,\mu$ on $\Omega_n$, use
\[
 \|\nu-\mu\|_{\TV}
 =\tfrac12\sum_x|\nu(x)-\mu(x)|
 =\sup_{A\subseteq\Omega_n}|\nu(A)-\mu(A)|.
\]
Define
\begin{align*}
 d_t^{\rm ct}(W,x)&=\|e^{t\mathcal L}(x,\cdot)-\mu_W\|_{\TV},&
 t_{\rm mix}&=\inf\{t\ge0:\max_xd_t^{\rm ct}(W,x)\le1/4\},\\
 d_k^{\rm dt}(W,x)&=\|P^k(x,\cdot)-\mu_W\|_{\TV},&
 k_{\rm mix}&=\min\{k\in\N\cup\{0\}:\max_xd_k^{\rm dt}(W,x)\le1/4\}.
\end{align*}
All limits in probability and statements holding with high probability refer to the disorder, together with explicitly introduced auxiliary randomness. No uniform assertion over disorder-adaptive time sequences is implicit.

We finish with the static identity that will turn overlap bounds into
covariance bounds. Zero field gives $\mu(x)=\mu(-x)$, so every spin has
mean zero. Consequently
\[
 \Sigma=\Cov_\mu(x)=\mu[xx^T],\qquad
 \Sigma_{ii}=1,
 \qquad q(x,y)=\frac{x^Ty}{n}.
\]
For independent samples from $\mu$, expansion of the square gives
\begin{equation}\label{eq:overlap-covariance}
 \Tr\Sigma^2=n^2\E_{\mu\otimes\mu}q^2.
\end{equation}
